\documentclass{article}
\usepackage{amsmath}
\title{L-System Fractal: Mathematical Foundation}
\author{Mathematica}
\begin{document}
\maketitle

\section{Introduction}

L-systems (Lindenmayer systems), introduced by Aristid Lindenmayer in
1968, are parallel rewriting grammars that model the self-similar growth
of biological structures. Iterating simple production rules generates
fractal geometry of arbitrary complexity.

\section{Formal Definition}

An L-system is a tuple $(\Sigma, \omega, P)$ where $\Sigma$ is the
alphabet, $\omega$ is the axiom, and $P$ is the production rule set.
After $n$ iterations, the string length grows as:
\[
  |S_n| = |S_{n-1}| \cdot \rho
\]
where $\rho$ is the branching factor. For the visualized system
$\rho = 2$, yielding $O(2^n)$ segments at iteration $n$.

\section{Hausdorff Dimension}

The fractal dimension $d_H$ of an L-system fractal relates its
self-similarity ratio $r$ to its copy count $N$:
\[
  d_H = \frac{\ln N}{\ln(1/r)}
\]
For the rendered 3D branching structure, $d_H \approx 1.585$, placing
it between a line and a plane --- a fractional dimension that captures
its intermediate complexity.

\section{Visualization Method}

Three iterations of the L-system axiom are evaluated into 3D
turtle-graphics coordinates. Each segment endpoint becomes a vertex in
the point cloud. The turntable camera animation reveals the branching
structure from all angles at 4K resolution.

\end{document}
